\documentclass[11pt]{article}

\usepackage[margin=1in]{geometry}
\usepackage{amsmath,amssymb,amsthm,mathtools}
\usepackage{bm}
\usepackage{xcolor}
\usepackage{hyperref}
\usepackage{enumitem}
\usepackage{microtype}

\hypersetup{colorlinks=true,linkcolor=blue!50!black,citecolor=blue!50!black,urlcolor=blue!50!black}

\newcommand{\R}{\mathbb{R}}
\newcommand{\Lap}{\nabla^2}
\newcommand{\Lapproj}{\nabla^2_{\!\mathrm{flat}}}
\newcommand{\Lapsphere}{\nabla^2_{\!\mathbb{S}}}
\newcommand{\Jflat}{\mathsf{J}_{\mathrm{flat}}}
\newcommand{\Jsphere}{\mathsf{J}_{\mathbb{S}}}
\newcommand{\meanS}[1]{\langle #1 \rangle_{\mathbb{S}}}
\newcommand{\Ld}{L_d}
\newcommand{\Ra}{\mathrm{Ra}}
\newcommand{\Pr}{\mathrm{Pr}}
\newcommand{\Ek}{\mathrm{Ek}}
\newcommand{\rjet}{r_{\mathrm{jet}}}
\newcommand{\phijet}{\varphi_{\mathrm{jet}}}
\newcommand{\bff}{\bm{f}}
\newcommand{\bu}{\bm{u}}
\newcommand{\bn}{\bm{n}}
\newcommand{\hr}{\hat{\bm{r}}}
\newcommand{\hphi}{\hat{\bm{\phi}}}
\newcommand{\hz}{\hat{\bm{z}}}

\title{Spherical NHQGE on a single polar cap via stereographic projection \\
{\large Formulation, spectral discretization, and discrete mean--fluctuation exchange}}
\author{Working note --- NHQGE / Jupiter polar dynamics}
\date{}

\begin{document}
\maketitle

\begin{abstract}
We extend the doubly-periodic Cartesian non-hydrostatic quasi-geostrophic
equation (NHQGE) solver to a spherical polar cap bounded poleward of a
midlatitude jet. Geometry is captured exactly via stereographic projection
from the opposite pole onto a tangent plane, which renders the induced
metric conformal with a single position-dependent factor.  The horizontal
operator that appears in the streamfunction-PV inversion is the spherical
Laplacian $\Lapsphere$ throughout, never the flat-disk Laplacian.  We
write down the prognostic equations on the cap, perform an azimuthal
Fourier decomposition, and discuss the radial spectral discretization in
detail: a Jacobi/Zernike basis (the recommended option, denoted
\emph{Option~A}) and a Chebyshev-on-$r^{2}$ alternative (\emph{Option~B}),
with the limitations of the latter near the pole made explicit.  The
vertical Chebyshev/CGL Galerkin/tau structure is unchanged from the
existing Cartesian code.  The dominant non-trivial change to the
mean--fluctuation thermal exchange is that all horizontal averages must
be taken with respect to the spherical area measure, not the projected
disk area; the discrete SBP corrector therefore operates on
spherical-area-weighted profiles, and the discrete adjoint pair must be
defined under the weighted inner product.  We give the time-stepping
algorithm, dealiasing prescription, boundary-condition enforcement,
diagnostic definitions, CFL constraint, and a validation strategy that
exploits the small-cap $\rjet \to 0$ Cartesian limit.
\end{abstract}

\tableofcontents

% -----------------------------------------------------------------------
\section{Domain and geometry}
% -----------------------------------------------------------------------

\subsection{Spherical cap}

Let $\mathbb{S}^2$ denote the unit two-sphere parameterized by latitude
$\varphi \in [-\pi/2, \pi/2]$ and longitude $\lambda \in [0, 2\pi)$.
Equivalently in colatitude $\vartheta = \pi/2 - \varphi \in [0, \pi]$,
$(\vartheta, \lambda)$ are standard spherical coordinates.  Throughout
we work in non-dimensional units (planet radius set to one); reinstating
a dimensional radius $a$ is a trivial rescaling.

The computational domain is a single polar cap from the north pole
($\vartheta = 0$) down to a fixed jet latitude $\phijet \in (0, \pi/2)$,
i.e.\ colatitude $\vartheta \in [0, \pi/2 - \phijet]$.  We treat the
boundary $\vartheta = \pi/2 - \phijet$ as a no-normal-flow wall (the
jet axis); see Section~\ref{sec:bcs} below.

\subsection{Stereographic projection from the south pole}

Project the sphere from the south pole onto the plane tangent to the
north pole.  A point at colatitude $\vartheta$ and longitude $\lambda$
maps to the projected disk coordinates
\begin{equation}
  r = \tan(\vartheta/2), \qquad \phi = \lambda,
  \label{eq:stereo-fwd}
\end{equation}
and the inverse is
\begin{equation}
  \vartheta = 2 \arctan(r), \qquad \lambda = \phi.
  \label{eq:stereo-inv}
\end{equation}
Under this map, the north pole maps to the origin $r=0$, the equator
maps to the unit circle $r = 1$, and the south pole maps to $r = \infty$.
The cap $\vartheta \in [0, \pi/2 - \phijet]$ maps to the disk
$r \in [0, \rjet]$ where
\begin{equation}
  \rjet \,=\, \tan\!\left(\frac{\pi/2 - \phijet}{2}\right)
        \,=\, \tan\!\left(\frac{\pi}{4} - \frac{\phijet}{2}\right).
  \label{eq:rjet}
\end{equation}
For $\phijet = 45^\circ$, $\rjet = \tan(22.5^\circ) \approx 0.4142$.
For $\phijet = 30^\circ$, $\rjet = \tan(30^\circ) = 1/\sqrt{3} \approx 0.5774$.

\subsection{Induced metric and conformal factor}

The line element on the sphere is $ds_{\mathbb{S}}^{2} = d\vartheta^{2}
+ \sin^{2}\vartheta \, d\lambda^{2}$.  Substituting~\eqref{eq:stereo-inv}
gives, after standard manipulation,
\begin{equation}
  ds_{\mathbb{S}}^{2} \;=\; \mu(r)\,\bigl(dr^{2} + r^{2} d\phi^{2}\bigr),
  \qquad
  \mu(r) \;=\; \frac{4}{(1+r^{2})^{2}}.
  \label{eq:metric}
\end{equation}
The map is therefore \emph{conformal}: the induced metric is the flat
metric scaled by a single position-dependent factor $\mu(r)$.  We will
use $\mu(r)$ throughout for this conformal factor.

The spherical area element is
\begin{equation}
  dA_{\mathbb{S}} \;=\; \mu(r)\,r\,dr\,d\phi.
  \label{eq:area}
\end{equation}
The total spherical area of the cap is
\begin{equation}
  A_{\mathrm{cap}}
  \;=\; \int_{0}^{2\pi}\!\!\int_{0}^{\rjet} \mu(r)\,r\,dr\,d\phi
  \;=\; 2\pi \int_{0}^{\rjet} \frac{4 r}{(1+r^{2})^{2}}\,dr
  \;=\; 2\pi\!\left(2 - \frac{2}{1+\rjet^{2}}\right)
  \;=\; \frac{4\pi\,\rjet^{2}}{1+\rjet^{2}}.
  \label{eq:Acap}
\end{equation}
For $\phijet = 45^\circ$: $A_{\mathrm{cap}} = 4\pi(\rjet^{2}/(1+\rjet^{2}))
\approx 4\pi \cdot 0.1716/1.1716 \approx 4\pi \cdot 0.1465$, i.e.\ about
$14.65\%$ of the full-sphere area $4\pi$. (This matches
$A_{\mathrm{cap}}/4\pi = (1 - \sin\phijet)/2$ for $\phijet = 45^\circ$,
giving $(1 - \sqrt{2}/2)/2 \approx 0.1464$, as expected.)

\subsection{Coriolis parameter and planetary vorticity gradient}

The Coriolis parameter on the rotating sphere is $f_{\mathbb{S}} =
2\Omega\sin\varphi = 2\Omega\cos\vartheta$.  Substituting
$\vartheta = 2\arctan(r)$,
\begin{equation}
  f(r) \;=\; 2\Omega\,\frac{1 - r^{2}}{1 + r^{2}}.
  \label{eq:fr}
\end{equation}
At the pole $f(0) = 2\Omega$; at the equator $f(1) = 0$; for the
$45^\circ$ cap edge, $f(\rjet) = 2\Omega(1-\tan^{2}(22.5^\circ))/
(1+\tan^{2}(22.5^\circ)) = 2\Omega \cos(45^\circ) = \sqrt{2}\,\Omega$.

The meridional gradient of $f$ along the sphere (the ``$\beta$ effect'')
is
\begin{equation}
  \beta_{\mathbb{S}}(r) \;\equiv\; \frac{\partial f}{\partial s_{\mathbb{S}}}
  \;=\; \frac{1}{\sqrt{\mu(r)}}\,\frac{df}{dr}
  \;=\; \frac{(1+r^{2})}{2}\,\frac{d}{dr}\!\!\left[2\Omega
  \frac{1-r^{2}}{1+r^{2}}\right]
  \;=\; -\,\frac{4\Omega\,r}{1+r^{2}},
  \label{eq:beta}
\end{equation}
where $s_{\mathbb{S}}$ is meridional arc length on the sphere and we used
$ds_{\mathbb{S}} = \sqrt{\mu}\,dr$.  Note that $\beta_{\mathbb{S}}(0) = 0$
(no $\beta$ at the pole) and $\beta_{\mathbb{S}}(\rjet) < 0$ in this
sign convention (i.e.\ $f$ decreases as one moves equatorward).

% -----------------------------------------------------------------------
\section{Differential operators on the sphere via the projection}
\label{sec:operators}
% -----------------------------------------------------------------------

For any 2D Riemannian metric $ds^{2} = g_{ij}\,dx^{i}dx^{j}$ with
inverse $g^{ij}$, the gradient and Laplacian of a scalar field $\psi$
are
\begin{equation}
  (\nabla \psi)^{i} \;=\; g^{ij}\,\partial_{j}\psi,
  \qquad
  \Lap \psi \;=\; \frac{1}{\sqrt{g}}\,\partial_{i}\!\left(
  \sqrt{g}\,g^{ij}\,\partial_{j}\psi\right).
  \label{eq:LaplaceBeltrami}
\end{equation}
For the conformally-flat metric~\eqref{eq:metric}, $g^{ij} =
\mu^{-1}\delta^{ij}$ and $\sqrt{g} = \mu r$, so the spherical
Laplacian and gradient in projected $(r, \phi)$ coordinates are
\begin{equation}
  \boxed{\;\;
  \Lapsphere \psi \;=\; \frac{1}{\mu(r)}\,\Lapproj \psi
  \;=\; \frac{1}{\mu(r)}\!\left[
  \frac{1}{r}\frac{\partial}{\partial r}\!\left(r\frac{\partial \psi}{\partial r}\right)
  + \frac{1}{r^{2}}\frac{\partial^{2}\psi}{\partial\phi^{2}}
  \right]
  \;\;}
  \label{eq:LapSphere}
\end{equation}
and
\begin{equation}
  \nabla_{\mathbb{S}} \psi
  \;=\; \frac{1}{\sqrt{\mu(r)}}\!\left(
  \hr\,\partial_{r}\psi + \hphi\,\frac{1}{r}\partial_{\phi}\psi\right).
  \label{eq:gradSphere}
\end{equation}
The two-dimensional Jacobian (the curl-projection of two gradients
along the surface normal $\hn$) is, equivalently,
\begin{equation}
  \boxed{\;\;
  \Jsphere(a, b)
  \;\equiv\; \nabla_{\mathbb{S}} a \times \nabla_{\mathbb{S}} b \cdot \hn
  \;=\; \frac{1}{\mu(r)}\,\Jflat(a, b),
  \;\;}
  \label{eq:JacSphere}
\end{equation}
with the flat-disk Jacobian
\begin{equation}
  \Jflat(a, b)
  \;=\; \frac{1}{r}\!\left[(\partial_{r} a)(\partial_{\phi} b)
  - (\partial_{\phi} a)(\partial_{r} b)\right].
  \label{eq:JacFlat}
\end{equation}
A scalar advection by an incompressible spherical-surface velocity
$\bu_{\mathbb{S}} = \hn \times \nabla_{\mathbb{S}}\psi$ is then
exactly $\bu_{\mathbb{S}} \cdot \nabla_{\mathbb{S}} a = \Jsphere(\psi, a)$,
and Eq.~\eqref{eq:JacSphere} states that the spherical advection
operator is the flat advection operator scaled by $\mu^{-1}$.

% -----------------------------------------------------------------------
\section{NHQGE on the sphere}
\label{sec:nhqge}
% -----------------------------------------------------------------------

\subsection{Continuum equations}

On the rotating sphere we keep the standard NHQGE prognostic structure
of four fields: perturbation potential vorticity $q'$, vertical
velocity $w$, perturbation temperature $\theta$, and mean-temperature
deviation $\overline{\Theta}'(z)$.  We non-dimensionalize as in the
existing Cartesian code (Miquel et al.~2026, eqs.~3.1a--d).  Replace
the flat horizontal operators by their spherical counterparts; the
vertical Chebyshev/CGL structure is unchanged.

\paragraph{Streamfunction--PV inversion.}
\begin{equation}
  \boxed{\;\;
  \bigl(\Lapsphere - \Ld^{-2}(z)\bigr)\,\psi(r,\phi,z)
  \;=\; -\,q'(r,\phi,z),
  \;\;}
  \label{eq:psi-q}
\end{equation}
to be solved \emph{pointwise in $z$} (the inversion does not couple
different $z$ levels).  The deformation radius $\Ld(z)$ may depend on
$z$ if stratification is non-uniform; for the Boussinesq NHQGE used
in the present code base $\Ld$ is constant in $z$ and the dependence
shown is for generality.

\paragraph{PV equation.}
\begin{equation}
  \partial_{t} q' + \Jsphere(\psi, q')
  + \Jsphere(\psi, f)
  - \partial_{z} w
  \;=\; -\,\mathcal{D}_{q}\,q',
  \label{eq:q-eq}
\end{equation}
with $\mathcal{D}_{q}$ a horizontal dissipation operator
(Section~\ref{sec:dissipation}).  Since $f = f(r)$ is axisymmetric, the
planetary stretching term simplifies to
\begin{equation}
  \Jsphere(\psi, f) \;=\; -\,\frac{1}{\mu(r)\,r}\,
  \frac{\partial \psi}{\partial \phi}\,\frac{df}{dr},
  \label{eq:Jpsi-f}
\end{equation}
which equals $\beta_{\mathbb{S}}(r)\,u_{r}^{\mathbb{S}}$ where
$u_{r}^{\mathbb{S}}$ is the meridional component of the spherical
physical velocity.

\paragraph{Vertical-velocity equation.}
\begin{equation}
  \partial_{t} w + \Jsphere(\psi, w)
  - \partial_{z}\psi
  - \frac{\Ra}{\Pr}\,\theta
  \;=\; -\,\mathcal{D}_{w}\,w.
  \label{eq:w-eq}
\end{equation}

\paragraph{Temperature-fluctuation equation.}
\begin{equation}
  \partial_{t}\theta + \Jsphere(\psi, \theta)
  + w\,\bigl(1 + \partial_{z}\overline{\Theta}'\bigr)
  \;=\; -\,\mathcal{D}_{\theta}\,\theta,
  \label{eq:theta-eq}
\end{equation}
where $\overline{\Theta}_{\mathrm{cond}}(z) = 1-z$ is the conductive
background and $\overline{\Theta}'(z)$ is its prognostic deviation.

\paragraph{Mean-temperature equation.}
\begin{equation}
  \boxed{\;\;
  \partial_{t}\overline{\Theta}'
  \;=\; -\,\partial_{z}\meanS{w\,\theta}(z)
  + \frac{1}{\Pr}\,\partial_{z}^{2}\overline{\Theta}',
  \;\;}
  \label{eq:meantemp-eq}
\end{equation}
where the bracket
\begin{equation}
  \meanS{F}(z) \;\equiv\;
  \frac{1}{A_{\mathrm{cap}}} \int_{0}^{2\pi}\!\!\int_{0}^{\rjet}
  F(r,\phi,z)\,\mu(r)\,r\,dr\,d\phi
  \label{eq:sphmean}
\end{equation}
is the \emph{spherical area mean} of $F$ over the cap.  The
distinction from the projected-disk mean (without the $\mu$ factor)
is the central numerical point of this document for the discrete
mean--fluctuation thermal exchange; see Section~\ref{sec:sbp}.

\subsection{Boundary conditions}
\label{sec:bcs}

\begin{description}[leftmargin=2em,style=nextline]
  \item[Vertical:] $w = \theta = \overline{\Theta}' = 0$ at $z = 0, 1$.
  $\psi$ inherits a horizontal-domain Helmholtz inversion at each $z$;
  no separate vertical BC on $q'$.
  \item[Pole, $r=0$:] regularity.  In azimuthal mode $m$, all fields
  must scale like $r^{|m|}$ near the origin to be smooth scalar
  functions of $(x,y)$.  This is built in to the Jacobi/Zernike basis
  (Section~\ref{sec:radialA}); for the Chebyshev-on-$r^{2}$ basis it
  has to be imposed by tau enforcement (Section~\ref{sec:radialB}).
  \item[Jet, $r=\rjet$:] no normal flow.  The radial component of the
  spherical physical velocity is
  $u_{r}^{\mathbb{S}} = -\,(1/\sqrt{\mu})(1/r)\,\partial_{\phi}\psi$,
  so $u_{r}^{\mathbb{S}}|_{r=\rjet} = 0$ for every $\phi$ requires
  $\partial_{\phi}\psi|_{r=\rjet} = 0$, i.e.\ $\psi(r=\rjet, \phi) =
  \psi_{\mathrm{jet}}(z, t)$ independent of $\phi$.  Without loss of
  generality choose the gauge $\psi_{\mathrm{jet}} \equiv 0$.  This
  becomes the Dirichlet condition on $\psi$ in the inversion
  Eq.~\eqref{eq:psi-q}, enforced by a tau row in the radial direction.
  We additionally impose Dirichlet $w = \theta = 0$ at $r = \rjet$
  for the dissipation operator and the Galerkin/tau treatment of the
  $w,\theta$ equations.
\end{description}

\subsection{Dissipation operators}
\label{sec:dissipation}

The Cartesian code uses operators of the form
$\mathcal{D}_{f} = \nu_{f}(-\Lap)^{p} + (\text{drag for }q')$
applied to each fluctuation field $f \in \{q', w, \theta\}$ with the
same hyper-order $p$.  On the sphere there are two natural choices
for the underlying horizontal Laplacian:

\begin{description}[leftmargin=2em,style=nextline]
  \item[Spherical-Laplacian dissipation:]
  $\mathcal{D}_{f} = \nu_{f}(-\Lapsphere)^{p}$.  Physically consistent.
  Costs more in the IMEX precompute because $\Lapsphere = \mu^{-1}
  \Lapproj$ is not diagonal in the radial basis: applying $\Lapsphere$
  in spectral space requires a multiplication by $\mu^{-1}$ in physical
  space (or a dense radial operator from the spectral form).
  \item[Flat-disk Laplacian dissipation:]
  $\mathcal{D}_{f} = \nu_{f}(-\Lapproj)^{p}$.  Cheaper and matches the
  Cartesian-code conventions exactly, but has the unphysical implication
  that dissipation is stronger near the pole (where $\mu$ is largest)
  than at the cap edge.  Acceptable as an experimental simplification
  if the dissipation amplitude is small enough to remain a sub-leading
  effect in the energy budget.
\end{description}

We adopt the spherical-Laplacian variant in this document and flag
where the cheaper variant would simplify the implementation.

% -----------------------------------------------------------------------
\section{Spectral expansion in azimuth and pole regularity}
% -----------------------------------------------------------------------

Expand any horizontal field $F(r,\phi,z,t)$ in azimuthal Fourier modes,
\begin{equation}
  F(r,\phi,z,t) \;=\; \sum_{m \in \mathbb{Z}} F_{m}(r,z,t)\,e^{i m \phi},
  \qquad
  F_{-m} = \overline{F_{m}} \text{ for real } F.
  \label{eq:Fourier-azim}
\end{equation}
For a real field we store only $m \ge 0$ via the standard real-FFT.
Linear horizontal operators with axisymmetric coefficients act
mode-by-mode.  In particular,
\begin{equation}
  \Lapproj F \;\longleftrightarrow\;
  \mathcal{L}_{m}\,F_{m}(r,z),
  \qquad
  \mathcal{L}_{m} \;\equiv\;
  \frac{1}{r}\frac{d}{dr}\!\left(r\frac{d}{dr}\right) - \frac{m^{2}}{r^{2}}
  = \frac{d^{2}}{dr^{2}} + \frac{1}{r}\frac{d}{dr} - \frac{m^{2}}{r^{2}},
  \label{eq:Lm-flat}
\end{equation}
and $\Lapsphere F \leftrightarrow \mu^{-1}(r)\,\mathcal{L}_{m}\,F_{m}$.
The Helmholtz inversion~\eqref{eq:psi-q} therefore reduces, in each
azimuthal mode $m$ and at each $z$, to a one-dimensional radial
problem
\begin{equation}
  \boxed{\;\;
  \mathcal{H}_{m}(z)\,\psi_{m}(r,z) \;=\; -\,q'_{m}(r,z),
  \quad
  \mathcal{H}_{m}(z) \;\equiv\;
  \mu^{-1}(r)\,\mathcal{L}_{m} - \Ld^{-2}(z),
  \;\;}
  \label{eq:Hm}
\end{equation}
with $\psi_{m}(\rjet, z) = 0$ as Dirichlet BC and the regularity
condition at $r = 0$ specified below.

\paragraph{Pole regularity.}
A smooth scalar function on the disk has a Taylor expansion
$F(r,\phi) = \sum_{j \ge 0}\sum_{m} c_{j,m} r^{j} e^{im\phi}$ in which
each $(j, m)$ pair satisfies $j \ge |m|$ and $j - |m|$ is even.  In
azimuthal mode $m$, the radial profile $F_{m}(r)$ therefore admits an
expansion $F_{m}(r) = r^{|m|}\,G_{m}(r^{2})$ with $G_{m}$ analytic at
$r=0$.  Equivalently, $F_{m}(r)$ has a definite parity $(-1)^{|m|}$
under $r \to -r$ extended through the origin, and its leading behavior
is $F_{m}(r) \sim r^{|m|}$.  This is the regularity constraint that
the radial basis must respect.

% -----------------------------------------------------------------------
\section{Radial discretization, Option A: Jacobi/Zernike basis}
\label{sec:radialA}
% -----------------------------------------------------------------------

\subsection{Definition}

Let $\xi = 2(r/\rjet)^{2} - 1 \in [-1, 1]$ for $r \in [0, \rjet]$, so
$\xi = -1$ at the pole and $\xi = +1$ at the jet.  For each azimuthal
mode $m$, expand the radial profile of any field as
\begin{equation}
  \boxed{\;\;
  F_{m}(r) \;=\; \left(\frac{r}{\rjet}\right)^{|m|}\,
  \sum_{n=0}^{N_{r}-1} \widetilde{F}_{m,n}\,
  P_{n}^{(0,\,|m|)}(\xi),
  \;\;}
  \label{eq:Jacobi-basis}
\end{equation}
where $P_{n}^{(\alpha,\beta)}$ are the Jacobi polynomials of degree $n$
and parameters $(\alpha, \beta)$.  This is, up to normalization, the
Zernike basis on the unit disk rescaled to radius $\rjet$.  The leading
factor $(r/\rjet)^{|m|}$ enforces the correct $r^{|m|}$ behavior at
the pole \emph{exactly}, so no auxiliary regularity tau row is needed.

The Gauss--Jacobi quadrature appropriate to weight $w(\xi) =
(1-\xi)^{0}(1+\xi)^{|m|}$ provides $N_{r}$ nodes $\xi_{j}^{(m)}$ and
weights $w_{j}^{(m)}$ such that
\begin{equation}
  \int_{-1}^{1} g(\xi)\,(1+\xi)^{|m|}\,d\xi
  \;\approx\; \sum_{j=1}^{N_{r}} w_{j}^{(m)}\,g(\xi_{j}^{(m)})
\end{equation}
is exact for polynomials $g$ of degree at most $2N_{r} - 1$.  This is
exactly the quadrature needed to evaluate inner products of the
Jacobi polynomials $P_{n}^{(0,|m|)}$ against themselves with the
right weight.  In terms of $r$, the change of variable
$d\xi = (4 r/\rjet^{2})\,dr$ and $r\,dr = (\rjet^{2}/4)\,d\xi$ gives
\begin{equation}
  \int_{0}^{\rjet} F_{m}(r)\,r\,dr
  \;=\; \frac{\rjet^{2}}{4}\!\left(\frac{1}{2}\right)^{|m|}
  \int_{-1}^{1} \widetilde{F}_{m}(\xi)\,(1+\xi)^{|m|}\,d\xi,
\end{equation}
where $\widetilde{F}_{m}(\xi) = \sum_{n} \widetilde{F}_{m,n}
P_{n}^{(0,|m|)}(\xi)$ is the Jacobi-coefficient form (the leading
$(r/\rjet)^{|m|}$ factor and $\xi$-substitution combine into the
$(1+\xi)^{|m|}$ weight).  This is the natural inner product for
\emph{flat}-disk integrals.  For \emph{spherical} integrals we add a
factor of $\mu(r)$ in physical space; see Section~\ref{sec:sbp}.

\subsection{Differentiation in the Jacobi basis}

Standard Jacobi-polynomial recurrence gives the derivative
\begin{equation}
  \frac{d}{d\xi} P_{n}^{(0,|m|)}(\xi)
  \;=\; \tfrac{1}{2}(n + |m| + 1)\,P_{n-1}^{(1,\,|m|+1)}(\xi),
\end{equation}
and a chain rule converts $d/d\xi \to d/dr$ via
$d\xi/dr = 4r/\rjet^{2}$.  More usefully, one computes the matrix
$\mathsf{D}_{r}^{(m)}$ that maps Jacobi coefficients of $F_{m}(r)$ to
the Jacobi coefficients of $\partial_{r}F_{m}(r)$ in the appropriate
$(\alpha,\beta)$ family for the derivative; this is a sparse
banded matrix.  The radial Laplacian operator $\mathcal{L}_{m}$ in
spectral form is then
\begin{equation}
  \mathsf{L}_{m}
  \;=\; \mathsf{D}_{r}^{(m)\,2}
  + \mathsf{R}^{-1}\,\mathsf{D}_{r}^{(m)}
  - m^{2}\,\mathsf{R}^{-2},
  \label{eq:Lm-spectral}
\end{equation}
where $\mathsf{R}$ is the diagonal multiplication-by-$r$ operator in
physical space, lifted to spectral via the basis transform.  In
practice we form $\mathsf{L}_{m}$ as a single $N_{r} \times N_{r}$
dense matrix per $m$, then $\mathsf{H}_{m}(z) = \mathsf{M}^{-1}\,
\mathsf{L}_{m} - \Ld^{-2}(z)\,\mathsf{I}$ where $\mathsf{M}$ is the
multiplication-by-$\mu(r)$ operator in spectral form.  $\mathsf{M}$
is dense (it is the Jacobi expansion of the function $\mu(r)$ acting
multiplicatively); $\mathsf{H}_{m}(z)$ is therefore dense.  Inversion
is by LU once per $m$ (and once per $z$ if $\Ld(z)$ varies).

\subsection{Boundary conditions in the Jacobi basis}

Pole regularity: enforced by construction; nothing to do.

Jet Dirichlet $\psi_{m}(\rjet, z) = 0$: the value at $\xi = +1$ is
$\sum_{n} \widetilde{\psi}_{m,n}\,P_{n}^{(0,|m|)}(1)$ with the
identity $P_{n}^{(0,|m|)}(1) = 1$, so the BC is the linear constraint
$\sum_{n} \widetilde{\psi}_{m,n} = 0$.  We enforce it by replacing the
last spectral row of $\mathsf{H}_{m}(z)$ by this constraint (a tau
row).  For Neumann BCs (used optionally on $w, \theta$), use
$\partial_{r} P_{n}^{(0,|m|)}|_{\xi=1}$ given by the standard Jacobi
identity.

\subsection{Why this is the recommended option}

\begin{itemize}[leftmargin=1.5em]
  \item \textbf{Exact pole regularity.}  The factor $(r/\rjet)^{|m|}$
  is built in.  No tau enforcement near $r=0$, no spurious modes from
  trying to represent a non-polynomial behavior with a polynomial
  basis.
  \item \textbf{Spectral convergence.}  For analytic radial profiles,
  the Jacobi expansion converges geometrically.  The Cartesian code
  enjoys the same property in $z$ via Chebyshev expansion, so the
  full $(m, n, k_{z})$ representation is uniformly spectral.
  \item \textbf{Natural quadrature for nonlinear products.}
  Gauss--Jacobi nodes provide exact integration of polynomial radial
  products up to degree $2N_{r} - 1$, so over-resolution by a factor
  of $\sim 3/2$ in the radial direction yields a clean dealiased
  Jacobian (Section~\ref{sec:dealias}).
  \item \textbf{Same algorithmic shape as the existing Cartesian
  code.}  Per-$m$ dense radial solves play the same role as
  per-shell scalar solves in the Cartesian Helmholtz; precomputed and
  cached IMEX inverses generalize directly.
\end{itemize}

% -----------------------------------------------------------------------
\section{Radial discretization, Option B: Chebyshev-on-$r^{2}$}
\label{sec:radialB}
% -----------------------------------------------------------------------

For comparison, define $\xi = 2(r/\rjet)^{2} - 1 \in [-1, 1]$ as
before but expand the radial profile of every azimuthal mode in
plain Chebyshev polynomials of $\xi$:
\begin{equation}
  F_{m}(r) \;=\; \sum_{n=0}^{N_{r}-1} \widehat{F}_{m,n}\,T_{n}(\xi).
  \label{eq:Cheb-r2-basis}
\end{equation}
This basis is identical for every $m$ and uses standard Chebyshev
machinery (DCT, three-term recurrence, well-behaved differentiation
matrices), which is the only reason to consider it: the existing
Cartesian code already has this machinery for the vertical direction.

\subsection{The pole-regularity defect}

The basis~\eqref{eq:Cheb-r2-basis} is polynomial in $r^{2}$ but not
in $r$.  This is benign for $m = 0$ and for any $m$ such that
$|m|$ is even, where the regularity requirement $F_{m}(r) \sim
r^{|m|}$ near the pole is automatically a polynomial in $r^{2}$.
For $|m|$ odd, however, the regularity requirement is
$F_{m}(r) \sim r \cdot (\text{analytic in } r^{2})$ near the pole,
which is \emph{not} polynomial in $\xi = 2(r/\rjet)^{2} - 1$.  Near
$\xi = -1$ this requires representing $\sqrt{(1+\xi)/2}$, whose
Chebyshev expansion in $\xi$ converges algebraically (like $1/n^{2}$)
rather than spectrally.  The accuracy near the pole for any odd-$|m|$
mode therefore degrades to algebraic.

A weakened workaround is to enforce, by tau rows, that
$F_{m}(r=0) = 0$ for $|m| \ge 1$ (and also that low-order $r$
derivatives vanish for $|m| \ge 2$, etc.), but this only enforces a
finite number of vanishing conditions, not the correct $r^{|m|}$
analytic behavior, and does not recover spectral convergence.

\subsection{When Option~B is acceptable}

\begin{itemize}[leftmargin=1.5em]
  \item As a \emph{prototype} to validate the rest of the
  infrastructure (dealiasing, IMEX, mean exchange) before committing
  to the more elaborate Jacobi machinery.  In a prototype run,
  algebraic convergence near the pole is a known limitation, not a
  blocker.
  \item For \emph{axisymmetric} runs ($m = 0$ only, no shear
  instabilities), where the regularity defect does not appear at all.
  \item For \emph{benchmark cross-checks} of the Jacobi
  implementation, where Option~B run at high $N_{r}$ should agree
  with Option~A run at moderate $N_{r}$ in the smooth interior.
\end{itemize}

\subsection{Boundary conditions in Option~B}

The Dirichlet jet condition is the simplest sum:
$\sum_{n} \widehat{F}_{m,n}\,T_{n}(1) = \sum_{n} \widehat{F}_{m,n}
= 0$.  Enforced as a tau row.  Pole conditions, when imposed at all,
use $T_{n}(-1) = (-1)^{n}$, $T_{n}'(-1) = (-1)^{n+1}\,n^{2}$, etc.

% -----------------------------------------------------------------------
\section{Vertical discretization (Chebyshev/CGL Galerkin/tau)}
% -----------------------------------------------------------------------

The vertical discretization is taken over verbatim from the existing
Cartesian solver and is summarised here for completeness.  Let
$\eta = 2z - 1 \in [-1, 1]$ for $z \in [0, 1]$, expand each prognostic
field in Chebyshev polynomials of $\eta$, and store fields as
Chebyshev coefficients (the Galerkin/tau formulation).  Vertical
derivatives use the coefficient-space recurrence $\mathsf{G}_{Z}$ and
$\mathsf{G}_{Z}^{2}$.  The vertical Vandermonde $\mathsf{V}$ and its
analytic DCT-I inverse $\mathsf{V}^{-1}$ effect the
coefficient/nodal transform.  Boundary conditions on $w, \theta,
\overline{\Theta}'$ are enforced via tau-projection matrices that
replace the last two Chebyshev coefficients of each field by the BC
constraints.  The Dirichlet Galerkin basis $\{-T_{n} + T_{n+2}\}$ is
used for $w$ and $\theta$ in the implicit solve; $q'$ retains the
full Chebyshev coefficient representation; $\overline{\Theta}'$ uses
the Dirichlet Galerkin basis on its $z$-only space.

The point of recapping this is that all the vertical infrastructure
ports unchanged: the spherical extension is purely horizontal.

% -----------------------------------------------------------------------
\section{Helmholtz inversion in spectral form}
\label{sec:helmholtz}
% -----------------------------------------------------------------------

In the Jacobi basis (Option~A), the inversion~\eqref{eq:Hm} for each
azimuthal mode $m$ becomes the dense linear system
\begin{equation}
  \boxed{\;\;
  \mathsf{H}_{m}^{(\mathrm{tau})}(z)\,
  \widetilde{\psi}_{m,\bullet}(z)
  \;=\; -\,\widetilde{q}'_{m,\bullet}(z),
  \;\;}
  \label{eq:Hm-system}
\end{equation}
where $\mathsf{H}_{m}^{(\mathrm{tau})}(z)$ is $\mathsf{H}_{m}(z) =
\mathsf{M}^{-1}\mathsf{L}_{m} - \Ld^{-2}(z)\mathsf{I}$ with the last
row replaced by the jet Dirichlet tau row, and the right-hand side has
its last entry zeroed out.

If $\Ld$ is independent of $z$ then $\mathsf{H}_{m}^{(\mathrm{tau})}$
depends only on $m$, and one factorization (LU) per $m$ suffices for
the entire run; there are $\sim N_{m} = N_{\phi}/2 + 1$ such
factorizations and the inversion at each $z$ level is one back-solve.
This is the analog of the Cartesian shell precomputation, with the
shells now being azimuthal modes rather than $|\bm{k}|^{2}$ shells.

If $\Ld(z)$ varies with $z$ (general stratification), there are
$N_{m} \times (N_{z}+1)$ factorizations, still tractable as a
once-per-run cost.

\paragraph{Cost.}  Per $m$, the factorization is $O(N_{r}^{3})$ once,
the application is $O(N_{r}^{2})$ per $z$ per step.  Total per step
is $O(N_{m} N_{z} N_{r}^{2})$ for the inversion alone, comparable to
the cost of the radial Jacobi transforms at each azimuthal mode.

% -----------------------------------------------------------------------
\section{Time stepping (ARS(2,2,2) IMEX)}
% -----------------------------------------------------------------------

We retain the ARS(2,2,2) IMEX scheme from the Cartesian solver:
diagonal coefficient $\gamma = 1 - 1/\sqrt{2} \approx 0.2929$,
off-diagonal $\delta = -\sqrt{2}/2$.  The implicit terms are the
linear vertical and horizontal dissipation, the buoyancy coupling
$\Ra/\Pr\,\theta$ in the $w$ equation, and the source $w$ in the
$\theta$ equation.  The explicit terms are the Jacobian advection,
the planetary-stretching $\Jsphere(\psi, f)$ term, and the
mean-exchange contributions in the thermal coupling
(Section~\ref{sec:sbp}).

\subsection{Implicit operator on the disk}

Per azimuthal mode $m$ and per Chebyshev coefficient $n_{z}$, the
implicit operator generalizes the Cartesian shell-by-shell scalar
construction to a $(N_{r}) \times (N_{r})$ dense matrix:
\begin{equation}
  \mathsf{A}_{m, n_{z}}^{\prime}
  \;=\; \alpha_{w,\mathrm{eff}}^{(m, n_{z})}\,\mathsf{I}
  - (\gamma \Delta t)^{2}\,\mathsf{C}_{m, n_{z}}\,\mathsf{B},
  \label{eq:Aprime}
\end{equation}
with
\begin{align}
  \alpha_{q}^{(m)}     &= 1 + \gamma \Delta t\,(\nu_{q}\,\mathsf{S}_{m,p} + \mathrm{drag}), \\
  \alpha_{w}^{(m)}     &= 1 + \gamma \Delta t\,\nu_{w}\,\mathsf{S}_{m,p}, \\
  \alpha_{\theta}^{(m)}&= 1 + \gamma \Delta t\,(\nu_{\theta}/\Pr)\,\mathsf{S}_{m,p}, \\
  \alpha_{w,\mathrm{eff}}^{(m,n_z)} &= \alpha_{w}^{(m)} - (\gamma \Delta t)^{2}\,(\Ra/\Pr)/\alpha_{\theta}^{(m)},
\end{align}
where $\mathsf{S}_{m, p}$ is the matrix representation of
$(-\Lapsphere)^{p}$ in the Jacobi basis at azimuthal mode $m$ (a
dense $N_{r} \times N_{r}$ matrix when the spherical Laplacian is
used; reduces to $\mathsf{L}_{m}$ for $p = 1$ scaled by
$\mu^{-1}$), $\mathsf{B} = \mathsf{G}_{Z}\,\mathsf{H}_{m}^{-1}\,
\mathsf{G}_{Z}$ is the vertical operator from the $\psi$-coupling
(now per $m$), and $\mathsf{C}_{m, n_{z}}$ is the corresponding
$\psi$-shell coefficient from the inversion.  Tau rows for the jet
Dirichlet condition replace the last spectral row of
$\mathsf{A}'_{m,n_{z}}$.

The block structure ($q', w, \theta$ jointly) is the same as in the
Cartesian code, with block elimination giving the
$\alpha_{w,\mathrm{eff}}$ form above.

\subsection{Precomputation and storage}

For Option~A with constant $\Ld$, $\mathsf{A}'_{m,n_{z}}$ depends on
the pair $(m, n_{z})$ and is precomputed once.  Number of distinct
matrices: $N_{m} \times (N_{z} + 1)$.  Storage: $N_{m} \times
(N_{z}+1) \times N_{r}^{2} \times 8$ bytes (dense, double).  For
$N_{m} = 64$, $N_{z}+1 = 65$, $N_{r} = 64$, this is $\approx
135\,\mathrm{MiB}$.  Comparable to the Cartesian IMEX storage.

% -----------------------------------------------------------------------
\section{Nonlinear Jacobian and dealiasing}
\label{sec:dealias}
% -----------------------------------------------------------------------

The Jacobian $\Jsphere(\psi, F) = \mu^{-1} \Jflat(\psi, F)$ is
evaluated pseudospectrally:
\begin{enumerate}[leftmargin=2em,itemsep=2pt]
  \item Transform $\psi, F$ from $(m, n)$ Jacobi/Fourier coefficients
  to a physical $(r_{j}, \phi_{k})$ grid.  Use Gauss--Jacobi nodes in
  $r$ (refined for dealiasing, see below) and uniform $\phi$ nodes
  with one inverse FFT per radial level.
  \item Compute $\partial_{r}\psi$, $\partial_{\phi}\psi$,
  $\partial_{r}F$, $\partial_{\phi}F$ in the same physical
  representation.  Radial derivatives are taken via the spectral
  derivative matrix $\mathsf{D}_{r}^{(m)}$ in coefficient space then
  inverse-transformed; azimuthal derivatives multiply the Fourier
  coefficient by $i m$ before inverse FFT.
  \item Form $\Jflat(\psi, F)$ pointwise on the $(r_{j}, \phi_{k})$
  grid via Eq.~\eqref{eq:JacFlat}.
  \item Multiply by $\mu^{-1}(r_{j})$ pointwise to obtain
  $\Jsphere(\psi, F)$.
  \item Forward-transform back to $(m, n)$ coefficients with the
  same Jacobi/Fourier weights.
\end{enumerate}

\paragraph{Azimuthal dealiasing.}  Standard $3/2$-rule on the
azimuthal FFT: pad from $N_{\phi}$ to $N_{\phi}^{\mathrm{pad}} =
3N_{\phi}/2$, compute the product, truncate back.  Identical to the
Cartesian code's horizontal dealiasing in one direction.

\paragraph{Radial dealiasing.}  The product of two functions of
azimuthal mode degree $|m_{1}|, |m_{2}|$ and radial polynomial
degree $n_{1}, n_{2}$ has azimuthal mode $m_{1} + m_{2}$ and radial
polynomial degree at most $n_{1} + n_{2} + (|m_{1}| + |m_{2}|)/2$
(the half-power coming from the $r^{|m|}$ prefactor when
$|m_{1}|, |m_{2}|$ have opposite parities).  The clean way to
dealias is to evaluate on a Gauss--Jacobi grid with
$N_{r}^{\mathrm{pad}} \ge \lceil 3N_{r}/2 \rceil$ nodes for the
\emph{worst-case} $|m| = N_{m} - 1$ Jacobi family, then truncate
back to $N_{r}$ coefficients in the appropriate per-$m$ Jacobi
family.  Mode-by-mode quadrature is exact for the relevant
polynomial degree when $N_{r}^{\mathrm{pad}}$ is chosen accordingly.

In Option~B (Chebyshev-on-$r^{2}$), the radial dealiasing is the
standard Chebyshev $3/2$-rule and is independent of $m$, but the
spectral content of the product near the pole is corrupted by the
$|m|$-odd half-power behavior, so radial dealiasing alone does not
restore accuracy.

\paragraph{Conservation properties.}  The Jacobian
$\Jsphere(\cdot,\cdot)$ is bilinear and antisymmetric, so the
identities
\begin{equation}
  \meanS{\psi\,\Jsphere(\psi, F)} = -\meanS{F\,\Jsphere(\psi, \psi)}/2 = 0
  \quad\text{(if }F=\psi\text{)},
\end{equation}
\begin{equation}
  \meanS{F\,\Jsphere(\psi, F)} = 0
  \quad\text{(any }F, \psi\text{)},
\end{equation}
should be preserved discretely up to truncation error.  The first
follows from antisymmetry $\Jsphere(\psi,\psi) = 0$; the second from
$\meanS{F\,\Jsphere(\psi, F)} = \tfrac{1}{2}\meanS{\Jsphere(\psi, F^{2})}$
plus the integration-by-parts identity
$\meanS{\Jsphere(\psi, G)} = 0$ for any $G$ when the BCs allow.
We discuss the integration-by-parts identity in detail below
because it is the discrete-conservation handle for the SBP corrector.

% -----------------------------------------------------------------------
\section{Integration-by-parts identity for the spherical Jacobian}
\label{sec:ibp}
% -----------------------------------------------------------------------

The cleanest statement of the conservation identity is:
\begin{lemma}
For any smooth scalar $G$ on the cap with the no-normal-flow
condition $\psi|_{r=\rjet} = 0$,
\begin{equation}
  \meanS{\Jsphere(\psi, G)} \;=\; 0.
  \label{eq:ibp}
\end{equation}
\end{lemma}
\begin{proof}
By Eq.~\eqref{eq:JacSphere} and~\eqref{eq:area},
\begin{align*}
  A_{\mathrm{cap}}\meanS{\Jsphere(\psi, G)}
  &= \int_{0}^{2\pi}\!\!\int_{0}^{\rjet}
  \mu^{-1}(r)\,\Jflat(\psi, G)\,\mu(r)\,r\,dr\,d\phi \\
  &= \int_{0}^{2\pi}\!\!\int_{0}^{\rjet}
  \Jflat(\psi, G)\,r\,dr\,d\phi \\
  &= \int_{0}^{2\pi}\!\!\int_{0}^{\rjet}
  \bigl[(\partial_{r}\psi)(\partial_{\phi}G)
  - (\partial_{\phi}\psi)(\partial_{r}G)\bigr]\,dr\,d\phi.
\end{align*}
The first term integrates to zero in $\phi$ (periodicity).  The
second term integrates to zero in $r$ via integration by parts: the
boundary contribution at $r=\rjet$ vanishes because $\psi(\rjet, \phi)
\equiv 0$, and at $r=0$ because $r\,\partial_{\phi}\psi \to 0$ for any
mode (regularity).  The interior derivative collapses by Leibniz to
$\partial_{r}(\partial_{\phi}\psi\cdot G)$ minus a term that
vanishes by periodicity in $\phi$.
\end{proof}

The crucial observation is that the conformal factor $\mu(r)$
\emph{cancels exactly} in the spherical area average of $\Jsphere$:
the $\mu^{-1}$ in the spherical Jacobian and the $\mu$ in the
spherical area element cancel.  Therefore the discrete identity to
preserve is the flat-disk identity $\int \Jflat(\psi, G)\,r\,dr\,d\phi
= 0$, which is independent of $\mu(r)$.  This is a clean discrete
target: any radial quadrature that integrates Jacobi-polynomial radial
products of high-enough degree exactly preserves Eq.~\eqref{eq:ibp}.

% -----------------------------------------------------------------------
\section{Mean--fluctuation thermal exchange under spherical area
weighting}
\label{sec:sbp}
% -----------------------------------------------------------------------

This is the dominant non-trivial change to the existing
\verb|balanced_sbp2_pc| machinery.  Read carefully.

\subsection{Continuum exchange identity}

Multiply Eq.~\eqref{eq:theta-eq} by $\theta$ and integrate over the
sphere using the area measure~\eqref{eq:area}:
\begin{equation}
  \tfrac{1}{2}\partial_{t} \meanS{\theta^{2}}
  + \meanS{\theta\,\Jsphere(\psi, \theta)}
  + \meanS{w\,\theta}\bigl(1 + \partial_{z}\overline{\Theta}'\bigr)
  \;=\; -\,\meanS{\theta\,\mathcal{D}_{\theta}\theta}.
  \label{eq:theta-energy}
\end{equation}
The advection term $\meanS{\theta\,\Jsphere(\psi,\theta)} =
\tfrac{1}{2}\meanS{\Jsphere(\psi, \theta^{2})} = 0$ by the IBP
identity~\eqref{eq:ibp}.  The remaining flux term
$\meanS{w\,\theta}\,\partial_{z}\overline{\Theta}'$ is exactly the
\emph{flux from fluctuation to mean} that must be balanced by the
\emph{flux from mean to fluctuation} in the
$\overline{\Theta}'$ equation.

Multiply Eq.~\eqref{eq:meantemp-eq} by $\overline{\Theta}'$ and
integrate over $z$:
\begin{equation}
  \tfrac{1}{2}\partial_{t}\!\int_{0}^{1}\!\overline{\Theta}'^{2}\,dz
  \;=\; -\!\int_{0}^{1}\overline{\Theta}'\,\partial_{z}\meanS{w\theta}\,dz
  + \frac{1}{\Pr}\!\int_{0}^{1}\overline{\Theta}'\,\partial_{z}^{2}\overline{\Theta}'\,dz.
  \label{eq:thbar-energy}
\end{equation}
Integration by parts in $z$ gives $-\int \overline{\Theta}'\,
\partial_{z}\meanS{w\theta}\,dz = \int \meanS{w\theta}\,\partial_{z}
\overline{\Theta}'\,dz$ (boundary terms vanish because
$\overline{\Theta}'|_{z=0,1} = 0$).  Adding this to the flux term in
Eq.~\eqref{eq:theta-energy}, the two flux contributions \emph{exactly
cancel}.  This is the continuum mean--fluctuation exchange identity.
The discrete machinery must reproduce this cancellation.

\subsection{Discrete spherical area mean}

In Option~A, the spherical area mean of any field $F(r, \phi, z)$ is
computed by:
\begin{enumerate}[leftmargin=2em,itemsep=2pt]
  \item Take the $m=0$ azimuthal Fourier coefficient $F_{0}(r, z)$
  (the $\phi$-average is $F_{0}$ exactly).
  \item Multiply by $\mu(r)$ and integrate against the radial measure
  $r\,dr$ from $0$ to $\rjet$ using Gauss--Jacobi quadrature with
  weight appropriate to the $m=0$ family:
  \begin{equation}
    \meanS{F}(z) \;=\; \frac{2\pi}{A_{\mathrm{cap}}}
    \int_{0}^{\rjet} F_{0}(r, z)\,\mu(r)\,r\,dr.
    \label{eq:disc-mean}
  \end{equation}
\end{enumerate}
The Gauss--Jacobi quadrature for $m=0$ uses weights $w_{j}^{(0)}$ at
nodes $\xi_{j}^{(0)}$, and the integrand is a polynomial in $r^{2}$
times $\mu(r) = 4(1+r^{2})^{-2}$.  Since $\mu$ is rational rather
than polynomial, the quadrature is not algebraically exact.  To
control the quadrature error, oversample radially when computing
mean-flux profiles $\meanS{w\theta}(z)$: use
$N_{r}^{\mathrm{mean}} \ge 3 N_{r}/2$ Gauss--Jacobi nodes for the
$m=0$ family, so that the polynomial part of the integrand is
integrated exactly to high order and the residual $\mu$
mis-quadrature is well-resolved.

\subsection{Discrete adjoint pair for the SBP corrector}

The current Cartesian SBP corrector implements the discrete pair
\begin{align}
  \text{(A)}\ \ &\text{fluctuation-side feedback to }\theta:
  \quad -\,w\,\partial_{z}^{(\mathrm{disc})}\overline{\Theta}', \\
  \text{(B)}\ \ &\text{mean-side flux exchange:}
  \quad -\,\partial_{z}^{(\mathrm{disc})}\overline{w\theta}^{(\mathrm{disc})},
\end{align}
where $\partial_{z}^{(\mathrm{disc})}$ is an SBP-2 difference matrix
on a uniform $z$-grid distinct from the CGL nodes, and the
$\overline{\,\cdot\,}^{(\mathrm{disc})}$ average is the Cartesian
horizontal area average.  The discrete adjoint requirement is that
\begin{equation}
  \sum_{j} \omega_{j}^{(\mathrm{sbp})}
  \overline{w\theta}^{(\mathrm{disc})}_{j}\,
  (\partial_{z}^{(\mathrm{disc})}\overline{\Theta}')_{j}
  + \sum_{j} \omega_{j}^{(\mathrm{sbp})}
  \overline{\Theta}'_{j}\,
  (\partial_{z}^{(\mathrm{disc})}\overline{w\theta}^{(\mathrm{disc})})_{j}
  \;=\; 0,
  \label{eq:disc-adjoint}
\end{equation}
for all admissible $\overline{w\theta}, \overline{\Theta}'$
satisfying the boundary conditions.  This is exactly the discrete
analog of the IBP identity in $z$ that produced cancellation in the
continuum.

\paragraph{Spherical change.}  Replace the Cartesian horizontal
average $\overline{\,\cdot\,}^{(\mathrm{disc})}$ everywhere by the
spherical area average $\meanS{\,\cdot\,}^{(\mathrm{disc})}$ defined
by Eq.~\eqref{eq:disc-mean}.  The $z$-direction adjoint structure of
Eq.~\eqref{eq:disc-adjoint} is unaffected and the SBP-2 difference
operator and weights $\omega_{j}^{(\mathrm{sbp})}$ on the uniform
$z$-grid remain identical to the Cartesian implementation.  What
changes is only the computation of the input profiles
$\meanS{w\theta}^{(\mathrm{disc})}(z)$ and the back-projection of
the mean flux from the SBP $z$-grid to the CGL solver $z$-grid (the
\verb|interp| transfer mode in the existing code).  The transfer
mode is a $z$-only operation and is independent of the horizontal
geometry.

\paragraph{Why this preserves global heat balance.}  The total heat
budget on the cap is
\begin{equation}
  \frac{d}{dt}\!\left[\tfrac{1}{2}\meanS{\theta^{2}}_{V}
  + \tfrac{1}{2}\!\int_{0}^{1}\overline{\Theta}'^{2}\,dz\right]
  \;=\; -\,\meanS{\theta\,\mathcal{D}_{\theta}\theta}_{V}
  + \frac{1}{\Pr}\!\int_{0}^{1}\overline{\Theta}'\,
  \partial_{z}^{2}\overline{\Theta}'\,dz,
\end{equation}
where $\meanS{\,\cdot\,}_{V}$ now includes a $z$-integration.  The
flux exchange terms cancel exactly in the continuum, and the SBP
discretization in $z$ preserves this cancellation \emph{by
construction}.  The horizontal averaging change (Cartesian flat
average $\to$ spherical area average) is purely a definition of what
``$\meanS{w\theta}(z)$'' means; once that average is computed
correctly, the SBP machinery in $z$ continues to produce a discrete
adjoint pair that closes the heat budget.

The danger we are explicitly avoiding: if one were to use the
flat-disk horizontal average for $\meanS{w\theta}(z)$ but then
update $\overline{\Theta}'(z)$ with the SBP corrector built around
that average, the $\overline{\Theta}'$ equation would be balancing
\emph{the wrong flux}, and the global heat budget would not close.
This would be the spherical analog of the dealiasing artifact that
plagued the early Cartesian runs.

\subsection{Exchange residual diagnostic}

Define
\begin{equation}
  \mathcal{R}_{\mathrm{exch}}(t)
  \;\equiv\; \int_{0}^{1}\!\!\!\meanS{w\theta}(z, t)\,
  \partial_{z}\overline{\Theta}'(z, t)\,dz
  \;-\; \int_{0}^{1}\!\overline{\Theta}'(z, t)\,
  \partial_{z}\meanS{w\theta}(z, t)\,dz,
  \label{eq:exch-resid}
\end{equation}
which the continuum identity sets to zero.  Its discrete realization
under the spherical area average and the SBP-2 $z$-derivative is the
diagnostic to monitor; the analog of \verb|mean_theta_exchange_resid\
ual_dealiased| in the Cartesian code.  An additional separate
diagnostic should report the discretization error of the spherical
area mean of $w\theta$ against the CGL-grid evaluation, analogous to
the \verb|Nusselt_dealiased| vs raw \verb|Nusselt| comparison.

% -----------------------------------------------------------------------
\section{Diagnostics under spherical area weighting}
\label{sec:diagnostics}
% -----------------------------------------------------------------------

All horizontal averages in diagnostics use Eq.~\eqref{eq:disc-mean}.
Specifically:

\paragraph{Nusselt number.}
\begin{equation}
  \mathrm{Nu}(t)
  \;=\; 1 + \int_{0}^{1}\!\!\meanS{w\theta}(z, t)\,dz
  \;-\; (\text{conduction reference}).
\end{equation}
With the closure $\overline{\Theta} = 1 - z + \overline{\Theta}'$,
\begin{equation}
  \mathrm{Nu}(t) \;=\; 1 + \int_{0}^{1}\meanS{w\theta}(z, t)\,dz.
\end{equation}

\paragraph{Kinetic-energy spectra.}  Spherical-area-weighted
horizontal KE shells are most naturally indexed by total wavenumber
$k_{n}^{(m)}$ associated with the $(m, n)$ Jacobi mode (the
eigenvalue of $-\Lapsphere$ on the corresponding eigenfunction in
the $\Ld \to \infty$ limit).  Define
\begin{equation}
  E_{h}(k, z, t)
  \;=\; \tfrac{1}{2}\sum_{(m, n) \in \mathrm{shell}_{k}}
  \,k_{n}^{(m)\,2}\,|\widetilde{\psi}_{m, n}(z, t)|^{2}\,N_{m,n},
\end{equation}
where $N_{m,n}$ is the squared norm of the $(m, n)$ Jacobi/Zernike
basis function under the spherical area inner product.  Vertical KE
spectra $E_{w}(k, z)$ are defined analogously using
$|\widetilde{w}_{m, n}|^{2}$ and the same shell binning.  The
Cartesian shell binning by $|\bm{k}|$ is replaced by a binning by
$k_{n}^{(m)}$ on the disk, but the post-processing pipeline is
otherwise unchanged.

\paragraph{Spectral KE flux.}  Define the cumulative KE flux
$\Pi(K) = -\sum_{k \le K} T(k)$ where $T(k)$ is the per-shell
transfer rate from the rotational and divergent contributions.  The
analog of the Cartesian $T_{h}(k) = \mathrm{Re}[\widehat{\psi}^{*}
\cdot \widehat{\Jflat(\psi, \zeta)}]$ becomes
$T_{h}(k) = \mathrm{Re}[\widetilde{\psi}^{*} \cdot
\widetilde{\Jsphere(\psi, \zeta)}]$ where the spectral inner product
uses the spherical area weight.  $T_{d}(k) = -\mathrm{Re}[
\widetilde{\psi}^{*} \cdot \widetilde{\partial_{z}w}]$ is unchanged
in form (it does not involve a horizontal Jacobian) but is summed
over Jacobi/Fourier shells under the spherical inner product.

% -----------------------------------------------------------------------
\section{CFL and resolution requirements}
% -----------------------------------------------------------------------

\paragraph{Horizontal CFL.}  Spherical CFL must use the spherical
arc-length metric.  In projected coords, the relevant horizontal
arc-length increment near radial node $r_{j}$ is $\Delta s_{j} =
\sqrt{\mu(r_{j})}\,\Delta r_{j}$, where $\Delta r_{j}$ is the local
radial node spacing.  The relevant velocity is the spherical
physical speed
$|\bu_{\mathbb{S}}| = \mu^{-1/2}\sqrt{(\partial_{r}\psi)^{2} +
r^{-2}(\partial_{\phi}\psi)^{2}}$.  The CFL constraint is then
\begin{equation}
  \Delta t \;\lesssim\; \min_{r, \phi, z}
  \frac{\sqrt{\mu(r_{j})}\,\Delta r_{j}}
  {|\bu_{\mathbb{S}}(r_{j}, \phi_{k}, z)|}.
\end{equation}
Note the $\mu^{-1/2}$ in the velocity and the $\sqrt{\mu}$ in the
arc length cancel into a factor of $\Delta r_{j}$ over the spectral
velocity, but the metric still affects the spacing $\Delta r_{j}$
itself: Gauss--Jacobi nodes cluster more or less near the endpoints
depending on the $|m|$ family.

\paragraph{Vertical CFL.}  Identical to Cartesian: the SBP corrector
imposes $\Delta t < \Delta z_{\mathrm{sbp}}/|w|_{\max}$ in the
explicit $w \partial_{z}$ step.  This is unaffected by the horizontal
geometry change.

\paragraph{Resolution.}  Spectral convergence in the Jacobi basis
requires $N_{r}$ proportional to the smallest physical horizontal
scale on the sphere, divided by the disk-to-sphere stretch factor at
that location: smaller scales near the cap edge require higher
$N_{r}$ than the same scales near the pole.  A practical rule of
thumb: pick $N_{r}$ such that the smallest resolved physical
wavelength on the sphere at $r=\rjet$ exceeds a few times the
expected dissipation scale.

% -----------------------------------------------------------------------
\section{Validation strategy}
% -----------------------------------------------------------------------

\subsection{Small-cap Cartesian limit}

For $\rjet \to 0$, $\mu(r) \to 4 + O(r^{2})$ uniformly across the
disk, $f(r) \to 2\Omega + O(r^{2})$, and the spherical Laplacian
reduces to $\Lapsphere \to (1/4)\Lapproj$.  Rescale
$\widetilde{r} = 2 r$, $\widetilde{\nabla} = 2\nabla$ to absorb the
factor of $4$, and the equations on the small cap reduce to the
existing Cartesian NHQGE on a disk of radius $2\rjet$ with $f =
2\Omega$ constant and $\beta = 0$.  A small-cap run should therefore
reproduce a known Cartesian-disk benchmark, providing an end-to-end
sanity check of the disk basis, IMEX, dealiasing, and mean exchange
machinery.

\subsection{Linear axisymmetric onset}

For the $m = 0$ subspace (purely axisymmetric), the linearized NHQGE
on the cap reduces to a generalized eigenvalue problem in the $m=0$
Jacobi radial basis.  Compute the lowest-Rayleigh-number unstable
eigenmode and compare to a pencil-and-paper computation of the same
eigenvalue (the integrated Sturm--Liouville problem in $r$ on
$[0, \rjet]$).  Cartesian analog: the existing $\Ra_{c} \approx
8.6956$, $k_{c} \approx 1.3048$ test for stress-free Rayleigh--B\'enard
in the doubly-periodic geometry.

\subsection{Discrete IBP and conservation tests}

\begin{itemize}[leftmargin=1.5em]
  \item Verify $\meanS{\Jsphere(\psi, F)}^{(\mathrm{disc})}$ is zero
  to roundoff for random $\psi$ satisfying $\psi|_{\rjet} = 0$ and
  random $F$, using both Option~A and Option~B radial bases.
  \item Verify the discrete exchange-residual~\eqref{eq:exch-resid}
  is zero to roundoff under randomly chosen $\meanS{w\theta}$ and
  $\overline{\Theta}'$ profiles.
  \item Verify that an integrated heat budget closes under linear
  forcing: $d/dt(\tfrac{1}{2}\meanS{\theta^{2}}_{V}) +
  d/dt(\tfrac{1}{2}\int \overline{\Theta}'^{2}) = -\text{(viscous
  drain)}$.
\end{itemize}

% -----------------------------------------------------------------------
\section{Algorithmic summary}
% -----------------------------------------------------------------------

The full algorithm in pseudo-code, per IMEX stage, with Option~A:

\begin{enumerate}[leftmargin=2em,itemsep=3pt]
  \item Given Jacobi/Fourier coefficients $\widetilde{q}'_{m, n}(z),
  \widetilde{w}_{m, n}(z), \widetilde{\theta}_{m, n}(z)$ at the
  current stage and Chebyshev coefficients of $\overline{\Theta}'(z)$:
  \item \textbf{Inversion.}  For each $m$, solve
  $\mathsf{H}_{m}^{(\mathrm{tau})}(z)\,\widetilde{\psi}_{m,\bullet}(z)
  = -\widetilde{q}'_{m,\bullet}(z)$ using the cached LU factor.
  Reconstruct $\partial_{z}\widetilde{\psi}_{m,\bullet}$ via the
  Chebyshev $\mathsf{G}_{Z}$ recurrence.
  \item \textbf{Transform to physical grid.}  Inverse Jacobi/Fourier
  transform $\widetilde{\psi}, \widetilde{q}', \widetilde{w},
  \widetilde{\theta}$ onto a $(r_{j}, \phi_{k})$ Gauss--Jacobi /
  uniform-$\phi$ grid, oversampled in both directions by $3/2$.
  Compute physical-space radial and azimuthal derivatives.
  \item \textbf{Jacobian evaluation.}  Compute
  $\Jflat(\psi, q'), \Jflat(\psi, w), \Jflat(\psi, \theta)$ on the
  oversampled physical grid.  Multiply by $\mu^{-1}(r_{j})$ pointwise
  to get $\Jsphere$.
  \item \textbf{Forward transform.}  Forward Jacobi/Fourier transform
  the spherical Jacobian fields back to truncated $(m, n)$
  coefficients.
  \item \textbf{Compute spherical horizontal means.}  For
  $F \in \{w\theta\}$, take $F_{0}(r_{j}, z)$ from the Fourier
  transform, and apply Eq.~\eqref{eq:disc-mean} to get
  $\meanS{w\theta}^{(\mathrm{disc})}(z)$.  Transfer to the SBP
  uniform-$z$ grid.
  \item \textbf{SBP corrector.}  Apply the
  \verb|balanced_sbp2_pc|-style discrete pair on the spherical
  area-weighted profiles to advance $\overline{\Theta}'$ and update
  the explicit thermal-exchange RHS for $\theta$.  Transfer back to
  the CGL solver $z$-grid.
  \item \textbf{Assemble explicit RHS.}  For each prognostic field,
  add Jacobian contributions, planetary-stretching contributions
  (from $\Jsphere(\psi, f)$), and mean-exchange contributions.
  \item \textbf{Implicit IMEX solve.}  For each $(m, n_{z})$ apply
  the cached $\mathsf{A}'_{m,n_{z}}{}^{-1}$ to the assembled implicit
  RHS to advance $(q', w, \theta)$ jointly to the next IMEX
  sub-stage.
  \item \textbf{Apply BCs.}  Tau-project for jet Dirichlet on $\psi$
  (already inside the inversion), and tau-project for vertical
  Dirichlet on $w, \theta, \overline{\Theta}'$.
  \item \textbf{Diagnostics.}  Periodically compute
  $\mathrm{Nu}(t)$, $E_{h}(k,z)$, $E_{w}(k,z)$,
  $\mathcal{R}_{\mathrm{exch}}(t)$, max horizontal/vertical CFL,
  etc.
\end{enumerate}

% -----------------------------------------------------------------------
\section{What is and is not changing relative to the Cartesian code}
% -----------------------------------------------------------------------

\paragraph{Unchanged.}  Vertical Galerkin/tau Chebyshev structure,
ARS(2,2,2) IMEX, block elimination of the buoyancy coupling,
\verb|balanced_sbp2_pc|-style discrete adjoint structure in $z$,
diagnostic categories (Nusselt, KE spectra, exchange residual),
overall outer solver loop and IO layout.

\paragraph{Changed.}  Horizontal Fourier $\to$ Jacobi/Zernike $\times$
azimuthal Fourier; flat horizontal Laplacian $\to$ spherical
Laplacian via $\mu^{-1}$ factor in the inversion; Cartesian Jacobian
$\to$ spherical Jacobian via $\mu^{-1}$ factor in evaluation;
horizontal mean for the SBP corrector $\to$ spherical area mean with
$\mu(r)\,r\,dr\,d\phi$ measure; Helmholtz inversion shell solve
$\to$ per-$m$ dense radial solve; KE-spectrum binning by $|\bm{k}|$
$\to$ binning by Jacobi/Fourier total wavenumber; CFL based on
spherical arc length, with the $\sqrt{\mu}\,\Delta r$ spacing.

\paragraph{New.}  Per-$m$ Gauss--Jacobi quadrature infrastructure;
basis-aware radial dealiasing; spherical-area-mean primitive used by
all diagnostics and the SBP corrector; jet Dirichlet tau row in the
radial Helmholtz inversion.

% -----------------------------------------------------------------------
\section{Open theoretical points to confirm with collaborators}
% -----------------------------------------------------------------------

These are the equation-side questions whose resolution is needed
before the implementation can claim to be a faithful spherical
NHQGE rather than a discretization of an unprincipled equation set.

\begin{enumerate}[leftmargin=2em,itemsep=3pt]
  \item \textbf{Confirm the QG operator is exactly $\Lapsphere$.}
  We assume throughout that the streamfunction--PV inversion is
  $(\Lapsphere - \Ld^{-2})\,\psi = -q'$ (one-and-a-half-layer QG on
  the sphere).  Confirm that the rapid-rotation asymptotic on the
  curved cap reduces the full primitive equations to this form
  without surviving curvature-correction terms (e.g.\ no extra
  $\propto \mu^{-2} f^{2}$ contributions in the inversion).
  \item \textbf{Confirm the planetary-stretching term form.}  We
  represent it as $\Jsphere(\psi, f)$, i.e.\ $f$ enters the PV
  equation purely through the Jacobian with the streamfunction.  This
  is what the materially-conserved-PV picture gives.  Confirm that
  no other $\beta$-like term appears (e.g.\ from the asymptotic
  expansion of horizontal divergence on a curved surface).
  \item \textbf{Confirm the temperature equation.}  The buoyancy
  $\Ra/\Pr\,\theta$ enters $w$ as in Cartesian; the term
  $w\,(1+\partial_{z}\overline{\Theta}')$ enters $\theta$ as in
  Cartesian.  No metric factor is expected because $z$ is the same
  vertical coordinate and the buoyancy is a vertical force.  Confirm
  this once on paper.
  \item \textbf{Stratification on the cap.}  $\Ld(z)$ may vary in $z$
  but is assumed independent of $r$ (no horizontal stratification).
  Confirm this is acceptable for the science target, or relax to
  $\Ld(r, z)$ (which the implementation supports as long as $\mathsf{H}_{m}$
  is recomputed for each $z$, with negligible additional cost).
\end{enumerate}

% -----------------------------------------------------------------------
\section{Summary}
% -----------------------------------------------------------------------

The spherical NHQGE on a single polar cap is implementable as a
direct generalization of the existing Cartesian solver provided
three points are respected throughout:

\begin{enumerate}[leftmargin=2em,itemsep=3pt]
  \item The QG operator in the $\psi$--$q'$ inversion is
  $\Lapsphere = \mu^{-1}\Lapproj$, and \emph{not} the projected-disk
  Laplacian.
  \item The horizontal mean entering the mean--fluctuation thermal
  exchange is the \emph{spherical} area mean
  $\meanS{\,\cdot\,}$ defined with the $\mu(r)$-weighted area
  measure, and the SBP-2 corrector operates on this average.  The
  flat-disk average will silently break the global heat budget.
  \item Pole regularity is built into the radial basis.  Option~A
  (Jacobi/Zernike with $r^{|m|}$ prefactor) does this exactly and
  is recommended.  Option~B (Chebyshev-on-$r^{2}$) is acceptable for
  axisymmetric prototypes and benchmark cross-checks but loses
  spectral accuracy near the pole for any odd-$|m|$ mode.
\end{enumerate}

The vertical Chebyshev/CGL Galerkin/tau structure, the
ARS(2,2,2) IMEX scheme, and the discrete-adjoint $z$-direction
machinery of the SBP corrector are inherited from the Cartesian code
unchanged.  All horizontal infrastructure (transforms, dealiasing,
implicit shell solves, diagnostics) is replaced by its disk-spectral
analog, with the conformal factor $\mu(r)$ entering through one
explicit pointwise multiplication in operator evaluations and one
explicit weight in horizontal averages.

\end{document}
